\documentclass[11pt]{article}

\usepackage[T1]{fontenc}
\usepackage[utf8]{inputenc}
\usepackage[english]{babel}
\usepackage{amsmath,amssymb,mathtools}
\usepackage{enumitem}
\usepackage[a4paper,margin=2.5cm]{geometry}
\usepackage[hidelinks]{hyperref}
\usepackage{microtype}

% -----------------------------------------------------------------------------
% Notation
% -----------------------------------------------------------------------------
\newcommand{\C}{\mathbb{C}}
\newcommand{\N}{\mathbb{N}}
\newcommand{\BH}{\mathcal{B}(H)}
\newcommand{\spann}{\operatorname{span}}
\newcommand{\rank}{\operatorname{rank}}
\newcommand{\id}{\mathrm{id}}

% Simple environments: deliberately without boxes or colored backgrounds.
\newenvironment{definition}{\par\medskip\noindent\textbf{Definition.}\ }{\par\medskip}
\newenvironment{example}{\par\medskip\noindent\textbf{Example.}\ }{\par\medskip}
\newenvironment{observation}{\par\medskip\noindent\textbf{Observation.}\ }{\par\medskip}
\newenvironment{exercise}{\par\medskip\noindent\textbf{Exercise.}\ }{\par\medskip}
\newenvironment{calculation}{\par\medskip\noindent\textbf{Calculation.}\ }{\par\medskip}

\title{October 19 Session\\Concrete Examples of $C^*$-Algebras}
\author{}
\date{}

\begin{document}
\maketitle



\section{Convention for the inner product}

Let $H$ be a complex Hilbert space and let $(e_i)$ be an orthonormal basis of $H$.
The convention used in the session is that the inner product is \emph{linear in
the first variable}. In particular,
\begin{align*}
\langle \alpha h_1+\beta h_2,g\rangle
 &=\alpha\langle h_1,g\rangle+\beta\langle h_2,g\rangle,\\
\langle h,h\rangle&\ge 0,
\qquad
\langle h,h\rangle=0\iff h=0,\\
\langle g,h\rangle&=\overline{\langle h,g\rangle}.
\end{align*}
Consequently, the second variable is conjugate-linear:
\[
\langle h,\alpha g_1+\beta g_2\rangle
=
\overline{\alpha}\,\langle h,g_1\rangle
+
\overline{\beta}\,\langle h,g_2\rangle.
\]

\section{Rank-one operators in $\mathcal{B}(H)$}

\begin{definition}
For $i,j$, define the operator
\[
 e_i\otimes e_j:H\longrightarrow H,
 \qquad
 (e_i\otimes e_j)(h):=\langle h,e_j\rangle e_i.
\]
In these notes the symbol $e_i\otimes e_j$ denotes this rank-one operator.
\end{definition}

Its range is
\[
(e_i\otimes e_j)(H)=\C e_i,
\]
and therefore
\[
\dim (e_i\otimes e_j)(H)=1.
\]
Thus, $e_i\otimes e_j$ is a rank-one operator.

Moreover, by the Cauchy--Schwarz inequality and because $\|e_i\|=\|e_j\|=1$,
\begin{align*}
\|(e_i\otimes e_j)(h)\|
 &=\|\langle h,e_j\rangle e_i\|\\
 &=|\langle h,e_j\rangle|\,\|e_i\|\\
 &\le \|h\|\,\|e_j\|\,\|e_i\|\\
 &=\|h\|.
\end{align*}
In particular,
\[
\|e_i\otimes e_j\|\le 1,
\qquad
 e_i\otimes e_j\in\BH.
\]

\subsection{Adjoint of the operators $e_i\otimes e_j$}

Write
\[
T_{ij}:=e_i\otimes e_j.
\]
Then
\[
T_{ij}^*=T_{ji},
\qquad\text{that is,}\qquad
(e_i\otimes e_j)^*=e_j\otimes e_i.
\]

Indeed, for $h,g\in H$,
\begin{align*}
\langle T_{ij}h,g\rangle
 &=\langle \langle h,e_j\rangle e_i,g\rangle\\
 &=\langle h,e_j\rangle\langle e_i,g\rangle\\
 &=\langle h,\langle g,e_i\rangle e_j\rangle\\
 &=\langle h,T_{ji}g\rangle.
\end{align*}
By the definition of the adjoint, it follows that $T_{ij}^*=T_{ji}$.

\subsection{Multiplication rule}

For indices $i,j,r,s$,
\[
(e_i\otimes e_j)(e_r\otimes e_s)
=
\delta_{jr}\,e_i\otimes e_s.
\]
Indeed, for $h\in H$,
\begin{align*}
(e_i\otimes e_j)(e_r\otimes e_s)(h)
 &= (e_i\otimes e_j)\bigl(\langle h,e_s\rangle e_r\bigr)\\
 &=\langle h,e_s\rangle\langle e_r,e_j\rangle e_i\\
 &=\delta_{jr}\langle h,e_s\rangle e_i.
\end{align*}

This is exactly the multiplication rule satisfied by matrix units.

\subsection{Composition with an arbitrary operator}

Let $T\in\BH$. In the photographs, the following identities are left as exercises:
\[
T\circ(e_i\otimes e_j)=T(e_i)\otimes e_j
\]
and
\[
(e_i\otimes e_j)\circ T=e_i\otimes T^*(e_j).
\]

The second identity follows directly from
\[
((e_i\otimes e_j)T)(h)
 =\langle Th,e_j\rangle e_i
 =\langle h,T^*e_j\rangle e_i.
\]

\section{Finite-rank projections}

For each $i$,
\[
p_i:=e_i\otimes e_i
\]
is a projection in $\BH$, since
\[
p_i^*=p_i
\]
and, using the multiplication rule,
\[
p_i^2=(e_i\otimes e_i)(e_i\otimes e_i)=e_i\otimes e_i=p_i.
\]
Moreover,
\[
p_i(H)=\C e_i,
\]
so $p_i$ has rank one.

\subsection{The corner determined by a rank-one projection}

Let $p_i=e_i\otimes e_i$ and let $T\in\BH$. Write
\[
T(e_i)=\sum_r \lambda_r e_r.
\]
Then the corner
\[
p_i\BH p_i
\]
consists of scalar multiples of $p_i$. In particular,
\[
p_i\BH p_i=\C p_i\cong\C,
\qquad
\lambda p_i\longmapsto \lambda.
\]

A direct way to identify the coefficient is
\begin{align*}
p_iTp_i(h)
 &=p_iT\bigl(\langle h,e_i\rangle e_i\bigr)\\
 &=\langle h,e_i\rangle p_iT(e_i)\\
 &=\langle h,e_i\rangle
   \left\langle \sum_r\lambda_r e_r,e_i\right\rangle e_i\\
 &=\lambda_i\langle h,e_i\rangle e_i\\
 &=\lambda_i p_i(h).
\end{align*}
Therefore,
\[
p_iTp_i=\lambda_i p_i.
\]

\begin{exercise}
If $A$ is a $C^*$-algebra and $p\in A$ is a projection, verify that
\[
pAp:=\{pap:a\in A\}
\]
is a $C^*$-subalgebra of $A$. In the photographs this subalgebra is called the
\emph{corner} determined by $p$.
\end{exercise}

In particular, when $p$ is a rank-one projection in a representation on a Hilbert
space,
\[
pAp\cong\C
\]
in the situation considered in class.

\subsection{Rank of a projection}

The terminology used in the session is
\[
p\text{ has rank one}
\quad\Longleftrightarrow\quad
\dim p(H)=1,
\]
and, more generally,
\[
p\text{ has rank }n
\quad\Longleftrightarrow\quad
\dim p(H)=n.
\]

\begin{example}
Consider
\[
p=e_1\otimes e_1+e_2\otimes e_2.
\]
Using the multiplication rule for matrix units,
\begin{align*}
p^2
 &= (e_1\otimes e_1+e_2\otimes e_2)^2\\
 &=e_1\otimes e_1+e_2\otimes e_2\\
 &=p.
\end{align*}
Moreover, $p^*=p$, and
\[
p(H)=\C e_1+\C e_2\cong\C^2.
\]
Therefore, $p$ is a rank-two projection.
\end{example}

\section{The corner $p\mathcal{B}(H)p$ for a rank-two projection}

Let
\[
p=e_1\otimes e_1+e_2\otimes e_2.
\]
For $T\in\BH$, write the coordinates of $T$ with respect to the orthonormal basis as
\[
T(e_r)=\sum_k\lambda_{kr}e_k.
\]
Then
\begin{align*}
(e_i\otimes e_j)T(e_r\otimes e_s)
 &= (e_i\otimes e_j)\bigl(T(e_r)\otimes e_s\bigr)\\
 &= (e_i\otimes e_j)
    \left(\sum_k\lambda_{kr}e_k\otimes e_s\right)\\
 &=\sum_k\lambda_{kr}(e_i\otimes e_j)(e_k\otimes e_s)\\
 &=\sum_k\lambda_{kr}\delta_{jk}\,e_i\otimes e_s\\
 &=\lambda_{jr}\,e_i\otimes e_s.
\end{align*}

Consequently,
\begin{align*}
pTp
 &= (e_1\otimes e_1+e_2\otimes e_2)
    T
    (e_1\otimes e_1+e_2\otimes e_2)\\
 &= (e_1\otimes e_1)T(e_1\otimes e_1)
  +(e_1\otimes e_1)T(e_2\otimes e_2)\\
 &\qquad
  +(e_2\otimes e_2)T(e_1\otimes e_1)
  +(e_2\otimes e_2)T(e_2\otimes e_2)\\
 &=\lambda_{11}e_1\otimes e_1
  +\lambda_{12}e_1\otimes e_2
  +\lambda_{21}e_2\otimes e_1
  +\lambda_{22}e_2\otimes e_2.
\end{align*}

Thus every element of $p\BH p$ is determined by four complex coefficients. As a
complex vector space,
\[
\dim_{\C}(p\BH p)=4.
\]
A natural basis is
\[
\{e_1\otimes e_1,\ e_1\otimes e_2,\ e_2\otimes e_1,\ e_2\otimes e_2\}.
\]
The multiplication rules
\[
(e_i\otimes e_j)(e_r\otimes e_s)=\delta_{jr}e_i\otimes e_s
\]
are precisely the matrix-unit multiplication rules of $M_2(\C)$.

\begin{exercise}
Let $S\in\BH$. Compute explicitly
\[
pSp.
\]
Then compute
\[
(pSp)(pTp).
\]
\end{exercise}

\section{The commutative example $\C^n$}

In the final part of the session, $\C^n$ is considered with coordinatewise
operations. For
\[
z=(z_i)_{i=1}^n,
\qquad
w=(w_i)_{i=1}^n,
\]
define
\[
zw=(z_iw_i)_{i=1}^n,
\qquad
z^*=(\overline{z_i})_{i=1}^n,
\]
and the norm
\[
\|z\|_\infty=\max_{1\le i\le n}|z_i|.
\]

With these operations, $\C^n$ is a unital commutative $C^*$-algebra with unit
\[
1=(1,\ldots,1).
\]
In the photographs this example is related to the representation theorem for
commutative $C^*$-algebras:
\[
\C^n\cong C(X),
\]
and the following question is left open:
\begin{center}
\emph{What is the space $X$ in this example?}
\end{center}



\end{document}
